\documentclass[tikz,border=10pt]{standalone}
\usepackage{tikz}
\usepackage{xcolor}

\definecolor{hse-blau75}{HTML}{8abde2}
\definecolor{hse-blau50}{HTML}{b4d3ed}
\definecolor{hse-blau25}{HTML}{dbe9f7}
\definecolor{hse-blau15}{HTML}{eaf2fa}
\definecolor{hse-hellgrau}{cmyk}{0,0,0,.08}

% Definiere die Parameter für die Dimensionen
\newcommand{\xdim}{20}  % x-Dimension (Breite)
\newcommand{\ydim}{20}  % y-Dimension (Höhe)
\newcommand{\transparencyFactor}{0.7}  % Transparenzfaktor

\begin{document}
\begin{tikzpicture}[scale=0.5]
    \pgfmathsetseed{2}  % Setze einen Seed für Reproduzierbarkeit
    
    % Erstelle ein Raster von Punkten mit zufälligen Verschiebungen
    \foreach \i in {0,...,\xdim} {
        \foreach \j in {0,...,\ydim} {
            \pgfmathsetmacro\dx{rand*0.5}
            \pgfmathsetmacro\dy{rand*0.5}
            \coordinate (p-\i-\j) at ({\i+\dx}, {\j+\dy});
        }
    }
    
    % Zeichne die Polygone
    \foreach \i in {0,...,\the\numexpr\xdim-1} {
        \foreach \j in {0,...,\the\numexpr\ydim-1} {
            % Prüfe, ob das Polygon innerhalb des Dreiecks liegt
            \pgfmathsetmacro\maxI{\xdim * (\j / \ydim)}
            \ifnum\i<\maxI
                \pgfmathsetmacro\rand{int(random(0,4))}
                \pgfmathsetmacro\xfactor{\i/\xdim}
                \pgfmathsetmacro\yfactor{\j/\ydim}
                % Setze die Transparenz so, dass sie nach unten und nach rechts zunimmt
                \pgfmathsetmacro\opacity{\transparencyFactor * (1 - \xfactor) * \yfactor}
                
                \ifcase\rand
                    \def\cellcolor{hse-blau75}
                \or
                    \def\cellcolor{hse-blau50}
                \or
                    \def\cellcolor{hse-blau25}
                \or
                    \def\cellcolor{hse-blau15}
                \or
                    \def\cellcolor{hse-hellgrau}
                \fi
                
                \pgfmathtruncatemacro\nextI{\i+1}
                \pgfmathtruncatemacro\nextJ{\j+1}
                
                \fill[\cellcolor, opacity=\opacity] 
                    (p-\i-\j) -- (p-\nextI-\j) -- (p-\nextI-\nextJ) -- (p-\i-\nextJ) -- cycle;
            \fi
        }
    }
\end{tikzpicture}
\end{document}
